\documentclass[11pt,a4paper]{article}
\usepackage[UTF8]{ctex}
\usepackage{geometry}
\usepackage{xcolor}
\usepackage{enumitem}
\usepackage{titlesec}
\usepackage{mdframed}
\usepackage{hyperref}
\usepackage{fontawesome5}
\usepackage{booktabs}
\usepackage{array}

\geometry{left=2.5cm, right=2.5cm, top=2.5cm, bottom=2.5cm}

% 颜色定义
\definecolor{interviewer}{RGB}{25,100,180}
\definecolor{candidate}{RGB}{22,140,90}
\definecolor{sectionblue}{RGB}{30,80,160}
\definecolor{notegray}{RGB}{100,100,100}
\definecolor{bglight}{RGB}{245,248,255}
\definecolor{bggreen}{RGB}{245,255,248}
\definecolor{bgorange}{RGB}{255,250,240}

% Q&A 环境
\newenvironment{QA}{%
  \begin{list}{}{%
    \setlength{\leftmargin}{0pt}%
    \setlength{\itemsep}{8pt}%
  }%
}{\end{list}}

\newcommand{\Q}[1]{%
  \item \begin{mdframed}[backgroundcolor=bglight, linecolor=interviewer, linewidth=1pt, innerleftmargin=10pt, innerrightmargin=10pt, innertopmargin=6pt, innerbottommargin=6pt]
  \textcolor{interviewer}{\textbf{\faComments\ 面试官：}}\ #1
  \end{mdframed}
}

\newcommand{\A}[1]{%
  \begin{mdframed}[backgroundcolor=bggreen, linecolor=candidate, linewidth=1pt, innerleftmargin=10pt, innerrightmargin=10pt, innertopmargin=6pt, innerbottommargin=6pt]
  \textcolor{candidate}{\textbf{\faUser\ 候选人参考回答：}}\ #1
  \end{mdframed}
}

\newcommand{\NOTE}[1]{%
  \begin{mdframed}[backgroundcolor=bgorange, linecolor=orange!60, linewidth=0.8pt, innerleftmargin=8pt, innerrightmargin=8pt, innertopmargin=4pt, innerbottommargin=4pt]
  \textcolor{orange!80!black}{\textbf{\faLightbulb\ 考察重点：}}\ \textcolor{notegray}{#1}
  \end{mdframed}
}

\titleformat{\section}{\Large\bfseries\color{sectionblue}}{}{0em}{}[\titlerule]
\titleformat{\subsection}{\large\bfseries\color{interviewer}}{}{0em}{}

\title{\textbf{模拟面试记录}\\
\large \textcolor{notegray}{Top Seed SeedDream / Canva 中国研究院 — 算法研究员岗位}\\
\normalsize \textcolor{notegray}{候选人：王冶 \quad 吉林大学人工智能学院博士生}}
\date{2026年4月}

\begin{document}

\maketitle
\tableofcontents
\newpage

%% =====================================================================
\section{开场与自我介绍}
%% =====================================================================

\begin{QA}
\Q{请用3分钟介绍一下你自己，重点讲你的研究方向和最具代表性的工作。}

\A{您好！我是王冶，吉林大学人工智能学院的博士生，预计2026年9月毕业。我的研究方向主要集中在\textbf{图像风格迁移、图像生成与编辑、多模态定制化生成}，以及\textbf{统一理解与生成}。

博士期间以第一作者发表了6篇论文，包括CVPR 2025、AAAI 2025两篇CCF-A类。最具代表性的工作是\textbf{OmniStyle}（CVPR 2025），我们构建了一个大规模可泛化的风格迁移框架，核心贡献在于设计了一套高质量数据过滤流水线，解决了以往风格迁移数据集质量差、规模小的瓶颈。

另一项代表性工作是\textbf{SigStyle}（AAAI 2025），面向个性化定制场景，用户只需提供少量参考风格图，系统就能准确捕获"签名式"风格，并迁移到任意内容图上。

在实习方面，我曾在阿里、昆仑万维、百度等公司做研究型实习，目前在腾讯优图参与青云计划，做统一理解与生成模型的工作。}

\NOTE{考察候选人对自身研究脉络的梳理能力，表达是否清晰、有重点，能否快速建立认知框架。}
\end{QA}

%% =====================================================================
\section{论文深度追问：OmniStyle（CVPR 2025）}
%% =====================================================================

\subsection{动机与问题定义}

\begin{QA}
\Q{OmniStyle的出发点是什么？在它之前，风格迁移领域存在哪些核心瓶颈？}

\A{风格迁移领域长期存在两个核心瓶颈：

\textbf{（1）数据质量问题：}以往的训练数据大多是用简单配对或爬取的方式获取的，风格图与内容图之间的配对质量很差——风格标签噪声大，内容与风格耦合不干净，导致模型难以学到"纯粹"的风格表示。

\textbf{（2）泛化能力有限：}以往方法在特定风格域上表现尚可，但面对新的、罕见的风格时泛化性严重下降。这背后根源仍然是训练数据质量和分布的问题。

OmniStyle的出发点就是从\textbf{数据}层面入手：如果我们能够大规模地过滤出真正高质量的风格迁移样本，就可以训练出泛化性更强的模型。}

\Q{具体说说你的数据过滤流水线（Filtering Pipeline）是怎么设计的？用了哪些判据？}

\A{我们设计了一个多阶段的过滤流水线，核心判据包括：

\textbf{（1）风格-内容解耦评估：}利用预训练的CLIP模型，分别提取图像的"内容特征"和"风格特征"，通过计算跨对之间的语义相似度，过滤掉内容与风格高度耦合的样本（例如，一幅梵高画，如果CLIP认为它和"星空"的内容相关性很强，那它就同时携带了太多内容信息，不适合作为纯风格参考）。

\textbf{（2）风格强度评估：}引入对抗性度量或基于Gram矩阵的风格强度评分，过滤掉风格不明显的样本。

\textbf{（3）内容完整性验证：}对内容图进行语义分割或对象检测，确保内容图不被风格污染。

通过这套流水线，我们从大规模爬取的数据中筛选出了高质量的训练集，这是OmniStyle泛化性能强的核心原因。}

\NOTE{重点考察候选人对数据工程的理解深度，是否能解释为什么数据质量是该任务的核心，以及具体过滤策略的技术细节。}

\Q{OmniStyle在模型架构上做了哪些设计？是基于扩散模型还是其他框架？}

\A{OmniStyle基于扩散模型框架，具体来说是在Stable Diffusion基础上扩展的。架构上的关键设计包括：

\textbf{（1）风格编码器：}我们设计了专门的风格编码器，将参考风格图编码为风格条件向量，注入到扩散模型的Cross-Attention层中，让去噪过程可以感知风格信息。

\textbf{（2）内容保持机制：}为了防止风格迁移过程中内容被破坏，我们引入了内容控制分支（类似ControlNet的思路），在保持内容结构的前提下完成风格化。

\textbf{（3）训练策略：}利用我们过滤好的高质量配对数据，在扩散模型上进行有监督的风格化训练，并设计了针对风格一致性和内容保真度的联合损失函数。}

\Q{OmniStyle的评估指标用的是什么？如何证明"风格迁移质量好"？}

\A{风格迁移的评估本身是一个难题，因为没有真实的ground truth。我们使用了以下指标：

\textbf{（1）风格相似度：}用Gram矩阵距离或CLIP风格特征相似度衡量输出图与风格参考图的风格一致性。

\textbf{（2）内容保真度（Content SSIM / LPIPS）：}用SSIM、LPIPS等指标衡量内容结构的保持程度。

\textbf{（3）用户研究（User Study）：}这是最有说服力的指标，我们做了大规模的人类偏好评测，让标注员对比不同方法的输出，OmniStyle在风格强度和内容保持的综合评分上显著优于baseline。

\textbf{（4）与现有SOTA对比：}和AdaIN、StyleFormer、DiffStyle等方法定量对比。}

\NOTE{考察候选人对评估方法论的掌握，尤其是在没有GT的生成任务中如何设计可信的量化评估。}
\end{QA}

%% =====================================================================
\section{论文深度追问：OmniStyle2（ECCV 2026 在投）}
%% =====================================================================

\begin{QA}
\Q{OmniStyle2的核心创新点是"去风格化驱动的风格迁移"（Destylize-Driven），这个思路是什么意思？为什么要这么做？}

\A{"去风格化"的核心思想是：与其直接学习"如何迁移风格"，不如先学习"如何去除风格"，即学习一个从风格化图像还原到内容图的逆过程。这样做的好处是：

\textbf{（1）获得更干净的内容表示：}通过去风格化，我们可以得到一个与风格无关的内容表示，这为后续的风格迁移提供了更纯净的内容基础。

\textbf{（2）数据自监督构建：}我们可以利用"风格化图像 $\rightarrow$ 内容图"这个方向作为自监督信号，构建双向训练范式。去风格化模型和风格化模型互相约束，形成循环一致性（Cycle Consistency）训练。

\textbf{（3）高保真输出：}由于内容提取更干净，最终风格化结果在内容保真度上显著提升，这正是OmniStyle2的核心优势。}

\Q{OmniStyle2与OmniStyle相比，主要解决了什么问题？两者的关系是什么？}

\A{OmniStyle解决了\textbf{泛化性}问题——通过大规模高质量数据训练，让模型能处理各种风格。但OmniStyle有一个局限：在某些复杂场景下，内容结构保持得不够精细，尤其是人脸、纹理细节等高频信息容易被风格"污染"。

OmniStyle2专门针对\textbf{高保真度}问题：在保持OmniStyle泛化性优势的基础上，通过去风格化-风格化的双向训练范式，显著提升了输出图像的细节保真度。可以说OmniStyle解决了"能不能迁移"，OmniStyle2解决了"迁移得好不好"。}

\NOTE{考察候选人对自己研究线的演进逻辑是否清晰，能否说清楚从OmniStyle到OmniStyle2的核心改进动机。}
\end{QA}

%% =====================================================================
\section{论文深度追问：SigStyle（AAAI 2025）}
%% =====================================================================

\begin{QA}
\Q{SigStyle和普通的风格迁移有什么本质区别？"Signature Style"怎么定义？}

\A{普通风格迁移是给定一张风格图，把另一张内容图转化为该风格。而SigStyle面向的是\textbf{个性化定制}场景：目标是捕获某个\textbf{特定艺术家或用户}的"标志性风格"（Signature Style），这种风格可能跨越多张参考图，具有一定的抽象性和一致性。

"Signature Style"的定义是：在多张参考作品中，超越单幅图像的表面纹理，能稳定复现的风格特征集合。比如梵高的"旋转笔触"、莫奈的"光影朦胧感"，这些是跨作品稳定存在的签名级风格。

技术上，SigStyle通过微调个性化文生图模型（类似DreamBooth / Textual Inversion）来学习这种签名风格，让模型将特定风格与一个文本token绑定，从而实现可文本控制的签名式风格迁移。}

\Q{SigStyle在训练时如何防止模型过拟合到参考图的内容上，而不是风格上？}

\A{这是个性化风格学习最核心的挑战——语义内容与风格的解耦。我们采取了几个关键措施：

\textbf{（1）数据增强策略：}对参考风格图进行大量内容级变换（随机裁剪、颜色抖动、物体替换等），让模型看到"不同内容但相同风格"的变体，强迫模型关注风格而非内容。

\textbf{（2）内容先验正则化：}在微调时加入内容正则化损失，约束模型在风格token上学到的不是内容信息。

\textbf{（3）文本解耦设计：}在Textual Inversion框架中，为风格设计专门的文本embedding，与内容文本相互独立，减少两者的特征干扰。

\textbf{（4）少步微调 + Early Stopping：}过多的微调步数会导致内容过拟合，我们通过验证集上的风格/内容解耦指标来确定最优停止点。}

\Q{SigStyle用户只需提供多少张参考图？模型的推理速度如何？有没有做few-shot甚至one-shot的实验？}

\A{SigStyle支持\textbf{少样本}场景，我们在论文中测试了1张到20张参考图的不同设置。实验表明：

\begin{itemize}
  \item \textbf{3--5张}参考图时，风格捕获质量已经相当好；
  \item \textbf{1张（one-shot）}时，性能有所下降，但对于具有强烈视觉特征的风格（如梵高）仍然可用；
  \item \textbf{10张以上}时，性能趋于饱和，继续增加参考图收益递减。
\end{itemize}

推理速度方面，由于底层是扩散模型，单张生成约需20--50步，在A100上大约1--3秒。我们也做了蒸馏实验，通过DDIM加速，可以将推理步骤压缩到5步以内，但质量略有损失。}

\NOTE{考察候选人对个性化生成技术栈（DreamBooth/Textual Inversion/LoRA）的掌握深度，以及对few-shot场景的理解。}
\end{QA}

%% =====================================================================
\section{论文深度追问：DP-Adapter（Graphical Models，CAD/Graphics 2025最佳论文提名）}
%% =====================================================================

\begin{QA}
\Q{DP-Adapter名字里的"Dual-Pathway"指的是什么两条路？各自的作用是什么？}

\A{DP-Adapter中的双路径（Dual-Pathway）指的是：

\textbf{Pathway 1 — 保真度路径（Fidelity Pathway）：}专注于捕获参考人物图像的外观细节，包括面部特征、服装纹理、肤色等，确保生成的人物与参考图高度一致。这条路径通过一个精细化的图像编码器提取高频视觉特征，并以AdaLN或Cross-Attention的方式注入扩散模型。

\textbf{Pathway 2 — 文本一致性路径（Text Consistency Pathway）：}专注于确保生成图像遵从文本描述，例如动作、场景、服饰颜色等可以被文本覆写的属性。这条路径增强了文本条件对图像生成的控制力，解决了单路径方法中"图像条件压制文本条件"的问题。

两条路径通过一个轻量级的交叉注意力融合模块协同工作，在保真度和文本可控性之间取得平衡。}

\Q{DP-Adapter为什么能获得CAD/Graphics 2025最佳论文提名？你认为它最有价值的贡献是什么？}

\A{我认为核心贡献在于精准定位了一个实际痛点：现有的人物定制化生成方法（如IP-Adapter、PhotoMaker等）普遍存在\textbf{"保真度与可控性的权衡困境"}——提高图像保真度往往会削弱文本控制能力，反之亦然。

DP-Adapter通过双路径设计，第一次明确地将这两个目标\textbf{解耦到两个独立的计算路径}中，而不是试图用单一特征空间同时满足两个相互竞争的约束。这个设计思路简洁但有效，在多个benchmark上相比IP-Adapter、InstantID等方法有显著提升。

此外，Adapter的轻量级设计使其可以插拔到现有的SD/SDXL等基础模型上，工程可落地性强，这也是评委认可的重要因素。}

\Q{与IP-Adapter、PhotoMaker、InstantID相比，DP-Adapter的优劣势各是什么？}

\A{
\textbf{对比IP-Adapter：}
\begin{itemize}
  \item 优势：DP-Adapter在人物保真度上更高，且文本可控性不因增强保真度而下降；
  \item 劣势：训练复杂度略高，需要同时优化两条路径的融合机制。
\end{itemize}

\textbf{对比PhotoMaker：}
\begin{itemize}
  \item 优势：DP-Adapter在单张参考图场景下表现更好；PhotoMaker更适合多图聚合场景；
  \item 劣势：多参考图聚合能力不如PhotoMaker。
\end{itemize}

\textbf{对比InstantID：}
\begin{itemize}
  \item 优势：DP-Adapter对服装、体型等非人脸属性的保持更好，不局限于人脸ID；
  \item 劣势：InstantID在人脸精细度上专项性更强，在人脸相似度指标上略胜。
\end{itemize}
}

\NOTE{考察候选人对相关工作的掌握程度，以及能否客观评价自己工作的优劣势——这是研究成熟度的重要体现。}
\end{QA}

%% =====================================================================
\section{技术深度：扩散模型基础}
%% =====================================================================

\begin{QA}
\Q{请从数学层面解释扩散模型的前向过程和反向过程，以及DDPM的训练目标是什么。}

\A{
\textbf{前向过程（Forward Process）：}
给定真实样本 $\mathbf{x}_0 \sim q(\mathbf{x}_0)$，前向过程逐步向数据中添加高斯噪声：
$$q(\mathbf{x}_t | \mathbf{x}_{t-1}) = \mathcal{N}(\mathbf{x}_t; \sqrt{1-\beta_t}\mathbf{x}_{t-1}, \beta_t \mathbf{I})$$
利用重参数化，可以直接采样任意时刻的噪声图：
$$\mathbf{x}_t = \sqrt{\bar{\alpha}_t}\mathbf{x}_0 + \sqrt{1-\bar{\alpha}_t}\boldsymbol{\epsilon}, \quad \boldsymbol{\epsilon} \sim \mathcal{N}(0, \mathbf{I})$$
其中 $\bar{\alpha}_t = \prod_{s=1}^{t}(1-\beta_s)$。

\textbf{反向过程（Reverse Process）：}
反向过程学习去噪分布 $p_\theta(\mathbf{x}_{t-1}|\mathbf{x}_t)$，目标是从纯噪声 $\mathbf{x}_T$ 还原数据。

\textbf{训练目标：}
DDPM的训练目标是最小化简化后的L2损失（噪声预测）：
$$\mathcal{L}_{simple} = \mathbb{E}_{t, \mathbf{x}_0, \boldsymbol{\epsilon}}\left[\|\boldsymbol{\epsilon} - \boldsymbol{\epsilon}_\theta(\mathbf{x}_t, t)\|^2\right]$$
即训练一个噪声预测网络 $\boldsymbol{\epsilon}_\theta$，让它准确预测第 $t$ 步加入的噪声，等价于对变分下界（ELBO）的简化优化。}

\Q{DDIM和DDPM的核心区别是什么？DDIM为什么可以用更少的步数推理？}

\A{
\textbf{核心区别：}
DDPM的前向过程是\textbf{Markov链}，每一步只依赖上一步；反向过程也依此设计，必须走完所有T步（通常T=1000）。

DDIM（Song et al., 2020）重新推导了一个\textbf{非Markov}的确定性前向过程，使得反向过程不再需要逐步随机采样，而是可以用\textbf{子序列}来跳步推理。

\textbf{为什么可以跳步：}
DDIM的反向更新公式为：
$$\mathbf{x}_{t-1} = \sqrt{\bar{\alpha}_{t-1}}\underbrace{\frac{\mathbf{x}_t - \sqrt{1-\bar{\alpha}_t}\boldsymbol{\epsilon}_\theta(\mathbf{x}_t,t)}{\sqrt{\bar{\alpha}_t}}}_{\text{predicted } \mathbf{x}_0} + \sqrt{1-\bar{\alpha}_{t-1}-\sigma_t^2}\boldsymbol{\epsilon}_\theta(\mathbf{x}_t,t) + \sigma_t\boldsymbol{\epsilon}$$
当 $\sigma_t=0$ 时，过程完全确定，可以选取任意步骤子集，20--50步即可获得与1000步相近的质量。}

\Q{Classifier-Free Guidance（CFG）的原理是什么？CFG scale 设得太高会有什么问题？}

\A{
\textbf{CFG原理：}
CFG（Ho \& Salimans, 2021）在推理时将条件得分和无条件得分线性外推：
$$\tilde{\boldsymbol{\epsilon}}_\theta(\mathbf{x}_t, c) = \boldsymbol{\epsilon}_\theta(\mathbf{x}_t, \emptyset) + w \cdot (\boldsymbol{\epsilon}_\theta(\mathbf{x}_t, c) - \boldsymbol{\epsilon}_\theta(\mathbf{x}_t, \emptyset))$$
其中 $w$ 是guidance scale，$c$ 是条件（如文本），$\emptyset$ 是空条件。训练时随机以概率 $p$ 丢弃条件（置为空），让同一模型学会有条件和无条件两种去噪。

\textbf{CFG scale过高的问题：}
\begin{itemize}
  \item \textbf{过饱和（Over-saturation）：}颜色极度饱和，图像看起来不自然；
  \item \textbf{过锐化（Over-sharpening）：}高频伪影增多，细节失真；
  \item \textbf{多样性下降：}生成结果趋于单一，缺乏随机性；
  \item \textbf{语义崩塌：}过度强调条件可能导致结构扭曲。
\end{itemize}
通常 $w \in [2, 10]$ 为合理范围，需要根据任务权衡。}

\Q{ControlNet的设计思路是什么？你在工作中是如何使用或借鉴ControlNet的？}

\A{
\textbf{ControlNet设计思路：}
ControlNet将预训练扩散模型的编码器部分复制一份作为可训练的"控制分支"，冻结原始模型，仅训练控制分支。控制分支接受额外的空间条件（如边缘图、深度图、姿态骨架等），通过零初始化的卷积层（Zero-Convolution）将控制信息注入到原模型的解码器中。

零初始化的意义：训练初始时，控制分支的输出为零，不影响原模型，随着训练进行，控制分支逐渐学会有效的控制信号，避免训练初期的不稳定。

\textbf{在我工作中的借鉴：}
在OmniStyle和DP-Adapter中，我借鉴了ControlNet"冻结主干+轻量可训练分支"的思想。对于内容保持，我用类似思路设计了内容控制分支；对于DP-Adapter，双路径的保真度路径也采用了类似的插拔式设计，使整体训练更稳定，且不破坏预训练模型的先验知识。}

\NOTE{这组问题考察候选人的扩散模型理论基础，是生成方向研究员必须掌握的核心知识。}
\end{QA}

%% =====================================================================
\section{技术深度：风格迁移专项}
%% =====================================================================

\begin{QA}
\Q{从AdaIN到扩散模型时代的风格迁移，技术发展经历了哪几个阶段？各阶段的核心思路是什么？}

\A{
\textbf{阶段一：优化类方法（Gatys et al., 2016）}
Gram矩阵作为风格表示，通过迭代优化使内容图的各层Gram矩阵逼近风格图，速度慢（每张图需数分钟），但开创了深度学习风格迁移的先河。

\textbf{阶段二：前馈网络类方法（Johnson et al., 2016）}
训练一个前馈网络替代迭代优化，速度大幅提升，但每个网络只对应一种风格，灵活性差。

\textbf{阶段三：自适应归一化类方法（AdaIN，Huang \& Belongie, 2017）}
将风格信息编码为特征的均值和方差，通过Adaptive Instance Normalization将内容特征的统计量对齐到风格特征，实现了任意风格的实时迁移。SANet、LinearTransfer等后续方法在此基础上增加了注意力机制。

\textbf{阶段四：GAN类方法（CycleGAN, StarGAN等）}
利用对抗训练和循环一致性，实现图像域间的无监督风格迁移，但往往需要大量特定域数据。

\textbf{阶段五：扩散模型时代（2022至今）}
InstructPix2Pix、StyleAligned、StyleCrafter等方法利用扩散模型强大的生成先验，配合Cross-Attention条件注入，实现了语义级的精细风格控制。这个阶段的方法在风格迁移质量和泛化性上均大幅超越前代。

OmniStyle即属于这个阶段，专注于解决扩散模型时代的数据质量和泛化性问题。}

\Q{Gram矩阵作为风格表示有什么局限性？扩散模型时代的风格表示是如何改进的？}

\A{
\textbf{Gram矩阵的局限性：}
\begin{enumerate}
  \item \textbf{丢失空间信息：}Gram矩阵是特征通道间的相关性矩阵，空间维度被平均掉，导致无法捕获空间结构性的风格（如笔触方向、图案的空间排列）；
  \item \textbf{全局统计量的局限：}对于局部风格变化较大的图像，全局Gram矩阵无法精确描述局部风格特征；
  \item \textbf{语义感知弱：}Gram矩阵是底层纹理统计量，难以捕获高层语义级别的风格（如构图方式、光影处理）。
\end{enumerate}

\textbf{扩散模型时代的改进：}
\begin{enumerate}
  \item \textbf{CLIP特征空间：}利用CLIP的语义特征作为风格表示，语义感知能力更强；
  \item \textbf{Cross-Attention注入：}将风格图的特征序列（patch tokens）直接参与内容去噪过程的注意力计算，保留了更丰富的空间和语义信息；
  \item \textbf{Q-Former / Perceiver Resampler：}用轻量级Transformer将风格图特征压缩为固定长度的token序列再注入，兼顾效率和表达力（如IP-Adapter的设计）。
\end{enumerate}}

\Q{你如何看待"风格"这个概念的可计算定义？在你的研究中，你是如何操作化定义风格的？}

\A{这是一个哲学与工程的交叉问题。

从哲学上看，"风格"涵盖了颜色调性、笔触纹理、构图方式、光影处理、语义主题偏好等多个层次，很难用单一指标完全描述。

在我的研究中，我从以下几个层次操作化定义风格：

\textbf{（1）低层次（Low-level）：}颜色直方图、Gram矩阵，描述纹理和颜色调性；

\textbf{（2）中层次（Mid-level）：}VGG等CNN的中间层特征统计量，捕获笔触、图案等视觉风格元素；

\textbf{（3）高层次（High-level）：}CLIP图像编码器的特征，捕获语义级别的风格（如"水彩感"、"写实感"）；

\textbf{（4）个性化层次（Signature-level）：}通过Textual Inversion或LoRA微调捕获特定艺术家的跨作品稳定风格特征，这是SigStyle的核心操作。

在评估时，我倾向于用多层次的指标组合，而不依赖单一指标，因为单一指标很容易被针对性优化而失去代表性。}

\NOTE{这组问题考察候选人对风格迁移领域的全局视野和对"风格"这一核心概念的思辨能力。}
\end{QA}

%% =====================================================================
\section{当前研究项目追问}
%% =====================================================================

\begin{QA}
\Q{你在做的物理强化文生图（Physical-Aware Image Generation）是什么方向？解决了什么问题？}

\A{当前的文生图模型（如SD、FLUX等）生成的图像在视觉上很美观，但常常违反物理规律，例如：光影方向不一致、物体阴影缺失或方向错误、液体流动不合理、重力方向混乱等。

物理强化文生图的目标是让生成图像满足基本的物理约束。我的做法是：

\textbf{（1）物理先验注入：}将物理模拟（如光线追踪的简化模型、重力场方向）作为条件信号，通过ControlNet风格的分支注入扩散模型。

\textbf{（2）物理一致性损失：}训练时设计针对光影一致性、遮挡关系的约束损失，让模型在去噪过程中感知物理规律。

\textbf{（3）数据构建：}收集或合成具有准确物理标注的训练数据（如在3D渲染引擎中生成带物理GT的图像对）。

这个工作目标投稿ICML 2026，目前还在推进实验阶段。}

\Q{SVG图像生成是怎么做的？为什么要用扩散模型来生成SVG，而不是直接用代码生成？}

\A{SVG（Scalable Vector Graphics）生成是一个有趣的挑战。直接让LLM写SVG代码的方法存在明显局限：LLM对空间坐标的感知很差，生成的SVG代码往往在视觉上混乱，且难以控制细节。

我们的方法是将SVG生成转化为图像生成问题，再映射回矢量域：

\textbf{（1）图像生成阶段：}先用扩散模型生成高质量的光栅图像（Raster Image）；

\textbf{（2）矢量化阶段：}利用可微分的矢量化工具（如DiffVG）将光栅图像转换为SVG路径表示，或用神经网络直接预测贝塞尔曲线参数；

\textbf{（3）端到端优化：}将矢量化过程与扩散模型生成联合优化，使SVG在视觉效果和矢量表示的简洁性之间取得平衡。

这个方向的核心挑战是：如何让生成的SVG在视觉上高质量，同时路径数量精简，适合实际应用（如图标设计、插画）。}

\NOTE{考察候选人对当前在研工作的技术理解深度，以及能否向面试官清晰讲解还未发表的研究方向。}
\end{QA}

%% =====================================================================
\section{实习经历追问}
%% =====================================================================

\begin{QA}
\Q{腾讯优图青云计划的"Unified Understanding and Generation Model"是做什么的？你在其中负责哪个模块？}

\A{统一理解与生成（Unified Understanding and Generation）是当前大模型的重要趋势——将图像理解（如图像描述、VQA、分割）和图像生成（文生图、图像编辑）统一在同一个模型架构中，避免维护多个独立系统。

代表工作有Emu、Show-o、Janus等，它们通常使用自回归或扩散模型作为统一backbone，将理解任务和生成任务转化为同一个token预测框架。

在腾讯优图，我目前主要负责：
\begin{itemize}
  \item 生成模块的风格控制能力接入——将我在OmniStyle系列中积累的风格控制技术适配到统一模型框架；
  \item 多模态条件融合机制的设计，确保理解侧的特征能有效指导生成侧。
\end{itemize}
由于实习时间较短（2026年3月至今），具体细节仍在探索阶段。}

\Q{在阿里AMAP做Physical-Aware Image Generation时，具体落地场景是什么？有没有遇到过数据或工程层面的挑战？}

\A{在阿里AMAP（高德地图）的应用场景主要是\textbf{地图素材生成}，具体包括：
\begin{itemize}
  \item 生成符合真实光照条件的街景图（不同时段、天气下的场景渲染）；
  \item 地图标注图标的风格化生成（保持物理合理性的矢量图标扩充）。
\end{itemize}

\textbf{数据挑战：}物理正确的训练数据极其稀缺，真实世界的图像往往没有物理GT（如精确的光源位置、遮挡关系）。我们部分解决方案是利用3D渲染引擎生成合成数据，然后做域适应（Domain Adaptation）来弥合合成-真实域差距。

\textbf{工程挑战：}扩散模型的推理速度在高德的实际业务中是瓶颈（需要大量图片实时或批量生成），我参与了推理加速的工作，包括量化（INT8/FP16）和蒸馏（减少推理步数至8步）。}

\Q{昆仑万维的多主体人物图像生成和你后来发表的PairHuman有什么关联？}

\A{有直接关联。在昆仑万维实习期间，我参与了多主体人物生成（即在一张图中同时生成多个人物，保持各自ID一致性）的算法研发，这让我深刻认识到了该任务的核心挑战：

\textbf{主体混淆问题（Subject Mixing）：}当图中有两个人物时，扩散模型很容易将两人的外貌特征混淆，生成"杂交"的人物。

这个认识直接促使我在后续学术研究中构建了\textbf{PairHuman数据集}（发表于Information Fusion，中科院一区）——一个专门针对双人定制化生成的高保真摄影数据集，为解决主体混淆问题提供了数据基础。实习中遇到的工程痛点转化成了有价值的学术贡献，这是我觉得产学结合很重要的一个体现。}

\NOTE{考察候选人能否将实习经历与学术产出有机串联，体现研究嗅觉和问题驱动的研究方式。}
\end{QA}

%% =====================================================================
\section{针对 SeedDream 的专项问题}
%% =====================================================================

\begin{QA}
\Q{SeedDream是字节跳动旗下的图像生成大模型，你如何评价当前头部文生图模型（FLUX、SD3、Imagen3、SeedDream）的技术路线差异？}

\A{
\textbf{Stable Diffusion 3 / FLUX（Black Forest Labs）：}
采用Diffusion Transformer（DiT）架构，将扩散过程建立在Transformer backbone上（而非UNet），配合Flow Matching训练范式（代替DDPM）。FLUX在文字渲染、细节保真度上取得了显著突破，是目前开源生态最强的文生图模型。

\textbf{Google Imagen 3：}
继续沿用级联扩散（Cascaded Diffusion）路线，强调文本语义对齐和多样性。在人体解剖学准确性上做了专项优化。

\textbf{SeedDream（字节）：}
据公开信息，SeedDream采用了大规模中文美学数据的专项训练，在亚洲审美、人物生成上有明显优势。技术上应该也是DiT架构，结合了字节内部大规模算力和数据优势。

\textbf{我的判断：}
当前文生图的竞争已从"能不能生成"转向"生成质量的细节"：文字渲染准确率、人体手部细节、光影物理合理性、文本精准遵循度。这些恰好是我在研究中也关注的方向，我认为SeedDream在中文语境和亚洲审美上的专项优化是非常有价值的差异化策略。}

\Q{如果你加入SeedDream团队，你认为风格控制方向最值得突破的技术问题是什么？}

\A{我认为有三个最值得突破的方向：

\textbf{（1）风格的精细粒度控制：}当前的风格迁移要么粒度太粗（整体色调迁移），要么过拟合到参考图的内容上。如何实现"只迁移光影风格，不迁移构图"，或"只迁移笔触风格，不迁移配色"这样的精细粒度控制，是一个尚未被很好解决的问题。

\textbf{（2）视频/多帧风格一致性：}随着视频生成的兴起（Sora、Wan等），如何保证多帧之间的风格时序一致性，是风格控制在视频域的核心挑战。

\textbf{（3）风格与版权的边界问题：}商业产品中，如何精准复现某种"风格感"而不侵犯特定艺术家的版权，需要在技术（风格抽象化程度的控制）和产品策略上同时下功夫。这也是Canva、SeedDream这类面向大众用户的产品必须认真对待的问题。}

\NOTE{考察候选人对行业格局的理解，以及研究视野是否与业界前沿对齐。}
\end{QA}

%% =====================================================================
\section{针对 Canva 中国研究院的专项问题}
%% =====================================================================

\begin{QA}
\Q{Canva是一个面向设计师和大众用户的工具型产品，你认为生成式AI在Canva场景下最核心的技术挑战是什么？与纯学术研究有何不同？}

\A{Canva场景的核心挑战在于\textbf{可控性、实用性和安全性}，而不仅仅是生成质量的极致追求。

\textbf{（1）可控性：}用户需要生成符合品牌规范的图像（特定颜色、字体、风格调性）。学术方法中大而全的模型往往难以被精准约束，如何实现"品牌级"的精细风格控制是核心挑战。

\textbf{（2）文字渲染：}设计场景中文字是核心元素，而目前几乎所有文生图模型在文字渲染上都存在明显缺陷。这是Canva场景中一个极高优先级的技术问题。

\textbf{（3）推理速度与成本：}Canva的用户规模极大，模型必须足够高效才能商业可行。学术论文中经常用A100跑几分钟的方法，在生产环境中是不可接受的。

\textbf{（4）内容安全：}面向大众用户，生成内容的安全审核（NSFW过滤、版权合规）必须嵌入整个流程。

与纯学术研究的最大区别：学术关注"新颖性"，工业关注"可靠性"和"可扩展性"。一个方法能在benchmark上SOTA，不等于它在生产环境中可用。}

\Q{你的PaperABC学术频道有1.3万粉丝、28万播放量，你如何看待技术传播能力对一个研究员的价值？}

\A{我认为技术传播能力在今天的AI时代比以往任何时候都更重要，原因有三：

\textbf{（1）知识迭代加速：}AI领域每周都有大量新论文涌现，能够快速提炼、传播核心ideas，对自身学习和团队知识共享都有价值。运营PaperABC强迫我在理解一篇论文时必须深入到能"讲清楚"的程度，这显著提升了我的论文阅读质量。

\textbf{（2）建立技术影响力：}Canva这类产品公司需要吸引顶尖人才，研究院的声誉和外部技术影响力很重要。有能力做技术传播的研究员，对团队形象建设有额外价值。

\textbf{（3）跨界沟通能力：}研究员需要与产品经理、工程师有效沟通，把复杂技术讲得让非专业人士理解，是一项重要的软实力。

当然，传播能力永远是辅助，核心还是研究质量本身。}

\NOTE{Canva作为工具型公司，对候选人的工程落地意识、用户视角和沟通能力有更高要求。}
\end{QA}

%% =====================================================================
\section{系统设计题}
%% =====================================================================

\begin{QA}
\Q{请设计一个支持百万用户规模的在线风格迁移服务，用户上传内容图和风格图，返回风格化结果。要求延迟P99 < 3秒，描述你的系统架构。}

\A{这是一个典型的AI推理服务设计题，需要在\textbf{模型质量、推理速度、成本}之间取得平衡。

\textbf{整体架构：}
\begin{enumerate}
  \item \textbf{接入层：}API Gateway + CDN，负责请求路由、限流、鉴权。静态资源（结果图片）走CDN分发。

  \item \textbf{任务队列：}异步任务队列（Kafka/RabbitMQ）解耦前端请求和后端推理。用户上传图片后立即返回task\_id，前端轮询或WebSocket推送结果。

  \item \textbf{推理集群：}
  \begin{itemize}
    \item 模型：蒸馏后的扩散模型（如8步DDIM），配合TensorRT/TorchScript优化；
    \item 混合精度（FP16/INT8量化）将单次推理速度压缩到1--2秒（A10G/T4 GPU）；
    \item GPU批处理（Dynamic Batching），合并多个请求的推理，提高GPU利用率。
  \end{itemize}

  \item \textbf{存储层：}结果图片存储到对象存储（OSS/S3），数据库仅存元数据（task状态、用户信息）。

  \item \textbf{Auto-scaling：}基于队列长度自动扩缩容，低峰期节省成本，高峰期快速弹出。

  \item \textbf{缓存策略：}对相同（内容图hash + 风格图hash）的请求缓存结果，热门风格+热门内容组合可节省大量重复计算。
\end{enumerate}

\textbf{P99 < 3秒的保障：}通过蒸馏+量化将单次推理降至1.5秒，队列等待时间通过弹性扩容控制在1秒以内，总链路P99控制在2.5--3秒范围内。}

\NOTE{考察候选人的工程系统思维，能否跳出纯算法视角，从全栈角度思考AI服务落地。}

\Q{你有一个风格迁移模型，在A/B测试中发现"风格相似度"指标提升了10\%，但用户满意度（点赞率）反而下降了5\%。你怎么分析和解决这个问题？}

\A{这个情况说明我们的\textbf{自动化指标与用户体验之间存在偏差}，这是AI产品中非常常见但容易被忽视的问题。

\textbf{分析步骤：}

\textbf{（1）指标分解：}不能只看聚合的点赞率，需要分析是哪类用户、哪类内容/风格组合的满意度下降了。例如：风格相似度提升是否主要在某些极端风格上？这些风格可能对普通用户来说过于"重"，反而破坏了内容的可读性。

\textbf{（2）用户研究定性分析：}收集用户反馈，理解他们不满意的具体原因。常见原因可能是：风格太强烈压制了内容（用户其实更在意内容保真度）；生成结果颜色过于饱和；风格迁移破坏了文字/人脸的可读性。

\textbf{（3）重新审视"风格相似度"指标本身：}该指标是否存在对抗性——模型通过过度风格化来拿到高分，但这对用户实际上是伤害？需要引入\textbf{内容保真度约束}作为联合优化目标。

\textbf{解决方向：}重新设计训练损失，加入内容保真度惩罚项；或者在推理时引入风格强度控制参数（如OmniStyle中的CFG scale思路），让用户/系统自适应调节风格强度。}

\NOTE{考察候选人的产品思维，能否超越单一指标，从用户视角审视技术决策。}
\end{QA}

%% =====================================================================
\section{行为问题与综合素质}
%% =====================================================================

\begin{QA}
\Q{博士阶段你最受挫折的一次科研经历是什么？你是如何处理的？}

\A{最受挫折的一次是OmniStyle投稿CVPR之前的Rebuttal阶段。审稿人提出了一个质疑：我们的数据过滤流水线是否真的是性能提升的关键，还是仅仅是因为训练数据规模增大导致的？这个问题直击我们工作的核心贡献。

挫折感来自两方面：一是这个消融实验在提交前确实没有做完整；二是时间非常紧迫，Rebuttal只有72小时。

我的处理方式是：立即和导师讨论，拆解问题，把消融实验分工并行跑。72小时内，我们完成了控制变量实验（同等数据量，有无过滤流水线的对比），用数据证明了过滤策略本身带来了显著提升，而不仅仅是数据量增加的效果。最终Rebuttal说服了大多数审稿人，论文接收。

这次经历让我意识到：在做每个核心设计选择时，就应该同步设计对应的消融实验，而不是等到审稿后再补。这是我后续工作中养成的习惯。}

\Q{你同时有多篇在投论文，还在腾讯实习，如何管理时间和优先级？}

\A{我用一个比较朴实但有效的方式：按截止日期驱动，而不是按重要性驱动。

具体来说：
\begin{itemize}
  \item 每周日晚上列出本周所有事项，标注截止日期和依赖关系；
  \item 实习工作和科研工作\textbf{物理隔离}：上班时间专注实习项目，晚上和周末专注学术工作。两者都做到位，靠的是不让它们互相干扰；
  \item 在投论文的大部分工作已经完成，维护成本（回复审稿、修改实验）是间歇性的，可以见缝插针处理；
  \item 对于新项目，我会评估清楚投入产出比再启动，避免同时开启太多半成品。
\end{itemize}

说实话，这种高强度的多线程工作是有代价的——睡眠时间和社交时间都被压缩了。但在博士最后阶段冲刺期，我认为这是值得的权衡。}

\Q{如果面试结果是两家公司（SeedDream和Canva）都给了offer，你会如何选择？}

\A{这个问题我想诚实地回答，而不是说"哪个offer更好就选哪个"这样的套话。

从我的研究方向来看：\textbf{两家的重合度都很高}，风格迁移、图像生成都是核心业务场景。

\textbf{倾向SeedDream的因素：}字节体量大、算力资源充沛，做大规模模型训练的条件更好；SeedDream作为独立图像生成产品线，研究自主度可能更高。

\textbf{倾向Canva的因素：}Canva是面向真实用户的工具型产品，研究成果能更快看到落地效果；中国研究院体量相对小，个人成长空间和对核心方向的影响力可能更大。

最终我会重点考察：\textbf{直接mentor是谁、团队的研究文化（是research-driven还是product-driven）、以及给到的研究课题是否与我的积累高度匹配}。Offer本身是其次的，能让我把最近几年积累的风格迁移深度发挥出来、产生影响力，才是最重要的标准。}

\NOTE{这类行为问题考察候选人的自我认知、价值观清晰度，以及是否真诚——套话式回答会给面试官留下不好印象。}
\end{QA}

%% =====================================================================
\section{反向提问：候选人问面试官}
%% =====================================================================

以下是候选人可以向面试官提出的高质量问题，体现研究深度和对公司的了解：

\begin{enumerate}
  \item \textbf{关于研究自主度：}"团队的研究课题是自上而下分配，还是研究员可以根据自己的判断主导方向？我在风格迁移方向有较深积累，这方面是否有对应的研究课题？"

  \item \textbf{关于工程落地：}"研究成果落地到产品的周期通常是多久？研究院和产品线的协作模式是怎样的？"

  \item \textbf{关于数据资源：}"在图像生成方向，团队能够获取到哪些规模级别的训练数据？这对我设计实验方案有重要影响。"

  \item \textbf{关于发论文：}"公司对研究员发表学术论文的支持和政策是怎样的？研究成果是否可以在适当保护IP之后对外发表？"

  \item \textbf{关于mentor：}"我的直属导师或合作最密切的研究员是谁？他/她目前最关注的技术方向是什么？"
\end{enumerate}

%% =====================================================================
\section{总结与准备建议}
%% =====================================================================

\begin{mdframed}[backgroundcolor=bgorange, linecolor=orange!50, linewidth=1pt]
\textbf{面试前重点准备清单：}

\begin{enumerate}
  \item \textbf{论文精读：}对每篇第一作者论文，准备：(1) 一句话总结；(2) 核心贡献点；(3) 与前人工作的区别；(4) 局限性和未来方向。

  \item \textbf{扩散模型理论：}确保能流畅推导DDPM/DDIM/CFG的数学公式，理解Flow Matching的基本思路。

  \item \textbf{代码准备：}准备一个可以当场演示的demo（如OmniStyle或SigStyle的推理代码），展示工程能力。

  \item \textbf{了解目标公司：}
  \begin{itemize}
    \item SeedDream：关注字节最新的图像/视频生成进展，了解Seaweed-APT、SeedEdit等产品；
    \item Canva：了解Canva Dream Lab、Magic Studio等AI功能，思考Canva场景的独特技术需求。
  \end{itemize}

  \item \textbf{行为问题准备：}用STAR法则（情境-任务-行动-结果）准备3--5个典型故事，涵盖：克服困难、多方协作、创新方法、失败经历。

  \item \textbf{英文准备：}两家公司可能有英文技术交流，准备能用英文介绍自己最代表性的两篇工作。
\end{enumerate}
\end{mdframed}

\vspace{1em}

\begin{mdframed}[backgroundcolor=bglight, linecolor=interviewer!50, linewidth=1pt]
\textbf{候选人优势总结：}
\begin{itemize}
  \item \faCheckCircle\ 研究产出密度高，博士期间2篇CCF-A，发表节奏快；
  \item \faCheckCircle\ 风格迁移方向纵深强，OmniStyle系列有完整研究线；
  \item \faCheckCircle\ 大厂实习经历多元，兼具学术与工程视角；
  \item \faCheckCircle\ PaperABC自媒体展示出技术传播和影响力建设能力；
  \item \faCheckCircle\ CAD/Graphics最佳论文提名是额外亮点。
\end{itemize}

\textbf{需要注意的潜在弱点：}
\begin{itemize}
  \item \faExclamationTriangle\ 部分工作是CCF-B/C，顶会数量以CVPR/AAAI为主，需强调质量而非数量；
  \item \faExclamationTriangle\ 对LLM/多模态大模型方向的涉猎相对较少，Canva如有此类需求需补充；
  \item \faExclamationTriangle\ 在投论文较多，需要在面试中自信地表达这些工作的进展和把握度。
\end{itemize}
\end{mdframed}

\end{document}
